\documentclass[10pt,a4paper]{amsart}
\usepackage[margin=19mm]{geometry}
\usepackage{amsmath,amssymb,mathtools}
\usepackage[scr=rsfs,cal=euler]{mathalfa}
\usepackage{xfrac}
\usepackage[unicode,hidelinks,bookmarks=false]{hyperref}
\newcommand{\ie}{i.e.\ }
\newcommand{\cf}{cf.\ }
\newcommand{\resp}{respectively\ }
\newcommand{\Subsection}{\subsection}
\providecommand{\nobreakdash}{\nobreak\hbox{-}}
\setlength{\emergencystretch}{2em}
\multlinegap0pt
\usepackage{amssymb}
\usepackage{mathtools}
\usepackage{enumerate}
\newtheorem{proposition}{Proposition}
\title{The modified boundary contact process:\\ invariant measures and critical extinction}
\author{C\'elio Terra}
\date{Source: arXiv:2503.09732v2 (CC BY 4.0)}
\begin{document}
\maketitle

\noindent\emph{Source and licence.} The modified boundary contact process:\\ invariant measures and critical extinction. Version arXiv:2503.09732v2 (CC BY 4.0), distributed by arXiv under the Creative Commons Attribution 4.0 International licence, http://creativecommons.org/licenses/by/4.0/, as recorded in the arXiv licence metadata for this exact version and shown on the abstract page https://arxiv.org/abs/2503.09732v2. Full licence notice: \texttt{licenses/arxiv-2503.09732-CC-BY-4.0.txt}.\par\smallskip

\noindent\textit{Editorial scope.} Counted unit: the article's proposition prop\_uniquenessmeasure (source0.tex lines 193--204). The article's complete original proof of it is delivered with it (source0.tex lines 240--328, one \texttt{proof} environment of 89 lines, which the article places away from the statement). No mathematics is written, repaired or re-derived: the statement and the proof are the article's own lines, in the article's own notation, and no retained line is substituted or re-flowed.  No further source-local context is needed: every cross-reference of every retained line resolves inside the two delivered spans.  Nothing was omitted from any retained span, so the package carries no omission record.  No retained line contains a citation, so the wrapper carries no bibliography and no bibliography entry.  No retained line contains a figure environment, an image-inclusion command or drawing code, so no image is delivered and the package declares no asset dependency.  Editorial formatting only, in addition: an AMS \texttt{amsart} preamble that restates the article's own declarations and macros used by the retained lines, namely the environments \texttt{proposition} on separate counters, together with the standard packages those lines need.  The retained environments print in this document as Proposition 1 in order of appearance, whereas the article prints the same items with the numbering of its own full document.  The complete native source of the article as distributed in the arXiv e-print for this exact version is bundled unchanged as \texttt{sources/174661/source0.tex}.
	\begin{proposition} \label{prop_uniquenessmeasure}
		For every cylinder set $C\subseteq\Sigma^{\odot}$,
		\begin{equation*}
			\lim_{t\to+\infty}
			\left|
			\mathbb{P}_{\lambda_i,\lambda_e}(\Psi\xi_t^A\in C)
			-
			\mathbb{P}_{\lambda_i,\lambda_e}(\Psi\xi_t^B\in C)
			\right|=0,
		\end{equation*}
		uniformly in $A,B\in\Sigma^{\odot}$.
	\end{proposition}

	\begin{proof}[{Proof of Proposition~\ref{prop_uniquenessmeasure}.}]
	Let $(\zeta_t)_t$ be the auxiliary process with infection rate $\lambda_e$ only from the right edge to the right, and infection rate $\lambda_i$ everywhere else. More precisely, infections from $\mathcal{R}\zeta_t$ to $\mathcal{R}\zeta_t+1$ occur at rate $\lambda_e$, while all other nearest-neighbor infections occur at rate $\lambda_i$. We denote the law of this process by $\widehat{\mathbb{P}}_{\lambda_i,\lambda_e}$.
			
	Let $\tau_0:=0$ and, for $k\ge1$, let $\tau_k$ be the first time after $\tau_{k-1}$ at which the copy of the auxiliary process started from the origin at time $\tau_{k-1}$ dies out. More explicitly,
		\[\tau_k:=\inf\{t>\tau_{k-1}:\zeta_{\tau_{k-1},t}^{0}=\emptyset\}.\]
	Take $A\in\Sigma^{\odot}$. We define a recentered version $\tilde{\zeta}^A$ recursively on the intervals $[\tau_{k-1},\tau_k)$ by
		\[\tilde{\zeta}^A_t=\zeta_{\tau_{k-1},t}^{\Psi \eta}+\mathcal{R}\eta,\qquad \eta=\tilde{\zeta}^A_{\tau_{k-1}}.\]
	In words, at each time $\tau_{k-1}$ we recenter the current configuration at its right edge, restart the auxiliary process from that recentered configuration using the same graphical construction, and then translate the resulting configuration back to the previous right-edge position. By translation invariance of the graphical construction, $(\Psi\tilde{\zeta}^A_t)_{t\ge0}$ has the same distribution as $(\Psi\zeta^A_t)_{t\ge0}$ under $\widehat{\mathbb{P}}_{\lambda_i,\lambda_e}$. Since $A\in\Sigma^{\odot}$ is semi-infinite to the left, $(\Psi\zeta^A_t)_{t\ge0}$ under $\widehat{\mathbb{P}}_{\lambda_i,\lambda_e}$ has the same distribution as $(\Psi\xi^A_t)_{t\ge0}$ under $\mathbb{P}_{\lambda_i,\lambda_e}$.
	
	Define $N(t)=\sum_{k=1}^{\infty}\mathbf{1}_{\{\tau_k\le t\}}$. Fix $L<+\infty$, $\kappa\in\Sigma^{\odot}$, and consider the cylinder $C=\{D\in\Sigma^{\odot}:D\cap[-L,0]=\kappa\cap[-L,0]\}$. Every finite-dimensional cylinder for the process seen from the right edge can be written in this form for some $L$ and $\kappa$.
	
	For every $A,B\in\Sigma^{\odot}$,
	\begin{equation} \label{eq_difference_cylinder_bound}
		\begin{aligned}
			&|\mathbb{P}_{\lambda_i,\lambda_e}(\Psi\xi_t^A\in C)
			-\mathbb{P}_{\lambda_i,\lambda_e}(\Psi\xi_t^B\in C)|\\
			&\quad=
			|\widehat{\mathbb{P}}_{\lambda_i,\lambda_e}(\Psi\tilde{\zeta}_t^A\in C)
			-\widehat{\mathbb{P}}_{\lambda_i,\lambda_e}(\Psi\tilde{\zeta}_t^B\in C)|\\
			&\quad\leq
			\widehat{\mathbb{P}}_{\lambda_i,\lambda_e}
			\big(\Psi\tilde{\zeta}_t^A\cap[-L,0]
			\neq
			\Psi\tilde{\zeta}_t^B\cap[-L,0]\big).
		\end{aligned}
	\end{equation}
	
	Let $s=\tau_{N(t)}$.  The only point which differs from the classical
	contact process is that the set of allowed arrows may depend on the
	current configuration.  Here this difficulty is removed by the
	one-boundary auxiliary construction: after the recentering time $s$,
	both $\widetilde\zeta^A$ and $\widetilde\zeta^B$ are run from right edge
	$0$ with the same graphical representation, and the descendants of the
	site $0$ use exactly the arrows of the auxiliary process
	$\zeta^0_{s,t}$.  Thus the usual one-dimensional no-crossing argument
	applies to these active paths.  In particular,
	\[
	\left\{
	\Psi\widetilde\zeta^A_t\cap[-L,0]
	\neq
	\Psi\widetilde\zeta^B_t\cap[-L,0]
	\right\}
	\subseteq
	\left\{
	\inf\Psi\zeta^0_{\tau_{N(t)},t}>-L
	\right\}.
	\]
	Indeed, on the complement of the event on the right, the descendants of
	$0$ at time $s$ separate the window $[-L,0]$ from all sites initially to
	the left of $0$.
	
	By definition of $N(t)$, $\zeta^0_{\tau_{N(t)},t}\neq\varnothing$.
	Therefore $\{\inf\Psi\zeta^0_{\tau_{N(t)},t}>-L\}$ is contained in
	$\{0<|\zeta^0_{\tau_{N(t)},t}|\le L\}$, and hence
	\begin{equation}
		\widehat{\mathbb P}_{\lambda_i,\lambda_e}
		\left(
		\Psi\widetilde\zeta^A_t\cap[-L,0]
		\neq
		\Psi\widetilde\zeta^B_t\cap[-L,0]
		\right)
		\le
		\widehat{\mathbb P}_{\lambda_i,\lambda_e}
		\left(
		0<|\zeta^0_{\tau_{N(t)},t}|\le L
		\right).
		\tag{3.2}
		\label{eq_coupling_error}
	\end{equation}
	We now justify that the right-hand side of~\eqref{eq_coupling_error} tends to zero. Fix $L<+\infty$. There exists $\delta_L>0$ such that, for every non-empty finite set $F$ with $|F|\leq L$, the auxiliary process started from $F$ dies out during the next unit time with probability at least $\delta_L$. Indeed, this occurs if all particles in $F$ recover before time $1$ and no infection arrow issued from them is used before their recovery.
	
	Let $\rho=\widehat{\mathbb{P}}_{\lambda_i,\lambda_e}(\zeta_t^0\neq\emptyset\text{ for all }t\geq0)>0$. By the Markov property,
	\[
	\delta_L\,
	\widehat{\mathbb{P}}_{\lambda_i,\lambda_e}
	\big(0<|\zeta_s^0|\leq L\big)
	\leq
	\widehat{\mathbb{P}}_{\lambda_i,\lambda_e}
	\big(\zeta_s^0\neq\emptyset,\ \zeta_{s+1}^0=\emptyset\big).
	\]
	The right-hand side tends to zero as $s\to+\infty$, since $\widehat{\mathbb{P}}_{\lambda_i,\lambda_e}(\zeta_s^0\neq\emptyset)$ decreases to $\rho$. Hence
	\[
	\widehat{\mathbb{P}}_{\lambda_i,\lambda_e}
	\big(|\zeta_s^0|\leq L\,\big|\,\zeta_s^0\neq\emptyset\big)
	\longrightarrow0
	\quad\text{as }s\to+\infty.
	\]
	Since, after the first surviving trial, $N(t)$ is constant and $t-\tau_{N(t)}\to+\infty$, the right-hand side of~\eqref{eq_coupling_error} tends to zero. The convergence is uniform in $A$ and $B$, and the proposition follows
	\end{proof}

\end{document}
